\chapter{Що обчислюють \nt{GLUT}-нейрони}
\label{sec:glutcomputer}

\section{Правила гри}

У розділі~\ref{sec:neurontypes} ми побачили, що переважна більшість нейронів мозку --- \nt{GLUT}: вони лише збуджують, їхні ваги лише додатні. Візьмімо скільки завгодно таких нейронів і запитаймо: що може обчислити така мережа, якщо дозволяти собі лише те, що буває в живому організмі?

А в живому організмі немає тактового генератора, який перемикав би всі елементи одночасно. Кожен нейрон працює сам по собі, в неперервному часі: імпульси приходять і йдуть коли випаде, з різними затримками. Є \emph{входи} --- чутливі нейрони, що спрацьовують від світла, звуку чи тиску, --- і \emph{виходи} --- нейрони, що керують м'язами. Між ними --- мережа, і ніхто не каже їй, коли починати й коли закінчувати обчислення. Для електронника це \emph{асинхронна} схема з елементами, затримки яких заздалегідь точно не відомі.

Модель нейрона візьмімо найпростішу з тих, що чесно працюють у неперервному часі. Напруга $V$ на мембрані нейрона в спокої дорівнює нулю. Кожен імпульс, що надійшов від нейрона $j$, стрибком піднімає її на $w_j\ge0$, після чого напруга сама стікає назад до нуля зі сталою часу $\tau$:
\begin{empheq}[box=\keybox]{equation}
\tau\,\frac{\dd V}{\dd t} \;=\; -V \;+\; \tau\sum_j w_j \sum_k \delta\bigl(t-t_j^{(k)}-d_j\bigr), \qquad w_j\ge 0 .
\label{eq:glutneuron}
\end{empheq}
Тут $t_j^{(k)}$ --- моменти імпульсів нейрона $j$, $d_j$ --- затримка на шляху від нього, $\delta$ --- дельта-функція, яка й означає миттєвий стрибок. Коли $V$ досягає порога $\theta$, нейрон видає імпульс, $V$ скидається в нуль, і приблизно мілісекунду нейрон не може спрацювати знову. Типові числа: $\tau$ --- від часток мілісекунди до двадцяти мілісекунд у різних нейронів, затримки --- від часток мілісекунди до десятків мілісекунд. Це \emph{інтегрувальний нейрон із витоком}: RC-коло з порогом. Це свідоме спрощення: у нейронів кори гілки входів самі видають місцеві імпульси~\cite{Smith2013}, і один нейрон поводиться радше як невелика багатошарова мережа (розділ~\ref{sec:synthesis}). Для питань цього розділу досить і одного RC-кола.

Головна відмінність від логічного вентиля --- витік: нейрон пам'ятає імпульс не вічно, а приблизно $\tau$, і додається лише те, що надійшло в межах цього вікна.

\section{АБО та І}

Почнімо з вентилів. Якщо один імпульс піднімає напругу вище порога, $w\ge\theta$, нейрон спрацьовує на кожен імпульс будь-якого входу. Це \emph{АБО}.

З І цікавіше. Нехай $w<\theta<2w$: одного імпульсу замало, потрібні два. Але тепер питання «чи прийшли обидва» не має сенсу, доки не сказано, \emph{у межах якого часу}. Нехай перший імпульс прийшов у момент $0$, а другий --- через час $\Delta$. На момент приходу другого від першого залишилося $we^{-\Delta/\tau}$, і нейрон спрацює, якщо $we^{-\Delta/\tau}+w\ge\theta$. Звідси вікно збігу:
\begin{empheq}[box=\keybox]{equation}
\Delta \;\le\; \Delta_{\max} \;=\; \tau\,\ln\frac{w}{\theta-w} .
\label{eq:window}
\end{empheq}
Це не логічне І, а \emph{детектор збігів}: І в часі (рис.~\ref{fig:coincidence}). При $\theta=1.5w$ вікно дорівнює $\tau\ln2\approx0.7\tau$. У нейрона кори з $\tau=10\ms$ це близько семи мілісекунд. У нейронів слухової системи, де $\tau$ менша за мілісекунду, вікно стискається до часток мілісекунди.

\begin{figure}[t]
\centering
\begin{tikzpicture}[>=Stealth,x=0.9cm,y=1.6cm]
\foreach \dx/\lab/\c [count=\i] in {0.5/{збіглися}/red!70!black, 2.6/{не збіглися}/blue!60!black}{
 \begin{scope}[xshift={(\i-1)*7.2cm}]
  \draw[->,gray!70!black] (0,0) -- (6.3,0) node[below left,black] {\small $t$};
  \draw[->,gray!70!black] (0,0) -- (0,1.75) node[left,black] {\small $V$};
  \draw[densely dashed,gray!60!black] (0,1.5) -- (6.0,1.5) node[right,black] {\small $\theta$};
  \draw[very thick,\c] (0,0) -- (1,0) -- (1,1)
      plot[domain=1:1+\dx,samples=30] (\x,{exp(-(\x-1)/1.5)});
  \pgfmathsetmacro{\v}{exp(-\dx/1.5)+1}
  \draw[very thick,\c] (1+\dx,{exp(-\dx/1.5)}) -- (1+\dx,{min(\v,1.5)});
  \ifdim\v pt<1.5pt
    \draw[very thick,\c] plot[domain=1+\dx:6,samples=40] (\x,{\v*exp(-(\x-1-\dx)/1.5)});
  \else
    \draw[very thick,\c] (1+\dx,1.5) -- (1+\dx,0) -- (6,0);
    \draw[->,thick,\c] (1+\dx,1.55) -- (1+\dx,1.72) node[above,font=\scriptsize] {імпульс};
  \fi
  \draw[->,thick] (1,-0.35) -- (1,-0.08); \draw[->,thick] (1+\dx,-0.35) -- (1+\dx,-0.08);
  \draw[decorate,decoration={brace,amplitude=3pt,mirror}] (1,-0.42) -- (1+\dx,-0.42) node[midway,below=3pt] {\scriptsize $\Delta$};
  \node[\c] at (3.5,-0.95) {\small \lab};
 \end{scope}}
\end{tikzpicture}
\caption{Детектор збігів, поріг $\theta=1.5w$. Ліворуч другий імпульс прийшов через $\Delta<\Delta_{\max}$, сума досягла порога, нейрон спрацював. Праворуч $\Delta>\Delta_{\max}$: від першого імпульсу майже нічого не лишилося, і нейрон змовчав.}
\label{fig:coincidence}
\end{figure}

Інший спосіб отримати І --- через тривалість, а не через точний час. Поки горить світло, чутливий нейрон видає імпульси часто й рівно; якщо порогові замало одного такого входу, а двох досить, нейрон працює, поки активні обидва. Так мережа обчислює за \emph{рівнями}, і точний час приходу імпульсів неважливий. Більша частина того, що далі, буде саме так.

\section{Чого зробити не можна}

Тепер спробуймо зробити НЕ: нейрон, який працює, коли вхід мовчить. Не вийде: без вхідних імпульсів у~\eqref{eq:glutneuron} немає жодного стрибка, і напруга стоїть на нулі, нижче порога. А мережа з багатьох нейронів? Теж ні, і це доводиться відразу для всіх мереж.

Найзручніше говорити про частоти. Нехай $r_i$ --- частота імпульсів нейрона $i$, $I_i$ --- приплив від входів, а $f_i$ --- його передавальна характеристика: як частота залежить від сумарного припливу. В інтегрувального нейрона ця характеристика неспадна: більший приплив --- частіші імпульси. Частоти в мережі, що дійшла рівноваги, задовольняють рівняння
\begin{equation}
r_i \;=\; f_i\Bigl(\sum_j w_{ij}\,r_j + I_i\Bigr), \qquad w_{ij}\ge 0 .
\label{eq:ratenet}
\end{equation}

\begin{keyblock}
\begin{proposition}[Монотонність]
\label{prop:monotone}
Нехай мережа~\eqref{eq:ratenet} складається лише з \nt{GLUT}-нейронів і починає з повної тиші. Якщо посилити входи, то усталена частота жодного нейрона не зменшиться. Усе, що обчислює така мережа, --- \emph{монотонні} функції входів.
\end{proposition}
\end{keyblock}

Доведення. Почнімо з тиші, $r^{(0)}=0$, і підставляймо частоти в праву частину~\eqref{eq:ratenet}: $r^{(k+1)}=F\bigl(r^{(k)}\bigr)$. Права частина $F$ не спадає за жодним аргументом, бо ваги невід'ємні, а $f_i$ не спадають. Оскільки $r^{(1)}\ge0=r^{(0)}$, за індукцією $r^{(k+1)}\ge r^{(k)}$: частоти лише зростають і приходять до найменшого розв'язку~\eqref{eq:ratenet}. Якщо входи $I$ замінити на більші $I'$, то на кожному кроці $r^{(k)}(I)\le r^{(k)}(I')$, і в границі теж. Справжня мережа приходить до рівноваги не кроками, а неперервно, але для неї справджується те саме: за додатних зв'язків нерівність між двома прогонами, виконана на початку, зберігається весь час.

Для окремих імпульсів твердження правдиве не буквально: зайвий вхідний імпульс може змусити нейрон спрацювати трохи раніше, і через мілісекунду несприйнятливості пропаде його наступний імпульс. Але ефект слабкий і ненадійний, обчислювати на ньому не можна. Для частот монотонність строга.

Наслідки ті самі, що в будь-якої схеми без інверторів. Немає НЕ. Немає виключного АБО: додавання другого входу не може погасити вихід. І найнеприємніше: \emph{активність не можна зупинити активністю}, можна лише прибрати те, що її підтримує.

\section{Два проводи: вмикання й вимикання}

Вихід із глухого кута підказує сам організм. У багатьох органах чуття подразник передається двома наборами нейронів: у зорі одні клітини відповідають на посилення світла, інші --- на послаблення~\cite{Schiller1992}. Назвімо їх \emph{вмикальними} та \emph{вимикальними}. Замість одного проводу $x$ мережа отримує пару $(x,\bar x)$, і відсутність, яку вона не вміє обчислити, їй повідомляє сам датчик.

З парою проводів НЕ безплатне: досить поміняти проводи місцями. І та АБО робляться кожне двома нейронами:
\begin{empheq}[box=\keybox]{equation}
\begin{aligned}
x\wedge y \;&\to\; \bigl(\;x\wedge y,\;\; \bar x\vee\bar y\;\bigr),\\
x\vee y \;&\to\; \bigl(\;x\vee y,\;\; \bar x\wedge\bar y\;\bigr).
\end{aligned}
\label{eq:dualrail}
\end{empheq}
Праворуч усі операції монотонні, тож твердження~\ref{prop:monotone} не порушено.

Для асинхронної схеми двопровідний код має й другу перевагу, навіть важливішу. Пара проводів буває в трьох станах: «так», «ні» і, коли обидва мовчать, \emph{«відповіді ще немає»}. Отримувач дізнається, що обчислення закінчено, не за годинником, а за самим сигналом: у кожній парі запрацював рівно один провід. Між подразниками датчики мовчать, і мережа повертається в тишу, готова до наступного разу. В асинхронній електроніці на цьому прийомі будують схеми, які працюють правильно за будь-яких затримок елементів~\cite{Sparso2001}.

\begin{keyblock}
\begin{proposition}[Асинхронне обчислення]
\label{prop:dualrail}
Якщо кожен вхід надходить парою проводів «вмикання---вимикання», то мережа з \nt{GLUT}-нейронів може обчислити будь-яку логічну функцію входів. Відповідь теж надходить парами проводів, а її готовність видно з того, що в кожній парі запрацював один провід. Годинник для цього не потрібен. Ціна --- приблизно вдвічі більше нейронів.
\end{proposition}
\end{keyblock}

Приклад --- виключне АБО: «змінився рівно один із двох подразників». Прямий провід відповіді --- $(x\wedge\bar y)\vee(\bar x\wedge y)$, зворотний --- $(x\wedge y)\vee(\bar x\wedge\bar y)$. Це шість нейронів, і жоден із них не потребує гальмування (рис.~\ref{fig:dualxor}).

\begin{figure}[t]
\centering
\begin{tikzpicture}[>=Stealth,x=1cm,y=1cm,
  nrn/.style={circle,draw,very thick,fill=blue!6,minimum size=0.75cm,inner sep=0pt,font=\small},
  inp/.style={font=\small}]
\node[inp] (X)  at (0,3.3) {$x$};
\node[inp] (Xb) at (0,2.2) {$\bar x$};
\node[inp] (Y)  at (0,1.1) {$y$};
\node[inp] (Yb) at (0,0)   {$\bar y$};
\node[nrn] (A1) at (3.2,3.3) {$2$};
\node[nrn] (A2) at (3.2,2.2) {$2$};
\node[nrn] (A3) at (3.2,1.1) {$2$};
\node[nrn] (A4) at (3.2,0)   {$2$};
\node[nrn] (O)  at (7.2,2.75) {$1$};
\node[nrn] (Ob) at (7.2,0.55) {$1$};
\foreach \s/\t in {X/A1,Yb/A1,Xb/A2,Y/A2,X/A3,Y/A3,Xb/A4,Yb/A4}{\draw[->,thick,blue!60!black] (\s) -- (\t);}
\draw[->,thick,blue!60!black] (A1) -- node[pos=0.45,above,sloped,font=\scriptsize,black] {$x\wedge\bar y$} (O);
\draw[->,thick,blue!60!black] (A2) -- node[pos=0.45,below,sloped,font=\scriptsize,black] {$\bar x\wedge y$} (O);
\draw[->,thick,blue!60!black] (A3) -- node[pos=0.45,above,sloped,font=\scriptsize,black] {$x\wedge y$} (Ob);
\draw[->,thick,blue!60!black] (A4) -- node[pos=0.45,below,sloped,font=\scriptsize,black] {$\bar x\wedge\bar y$} (Ob);
\draw[->,very thick,red!70!black] (O) -- (8.6,2.75) node[right,black] {$x\oplus y$: «так»};
\draw[->,very thick,red!70!black] (Ob) -- (8.6,0.55) node[right,black] {$\overline{x\oplus y}$: «ні»};
\node[below,font=\small,gray!40!black] at (3.2,-0.55) {І: поріг $2$};
\node[below,font=\small,gray!40!black] at (7.2,-0.55) {АБО: поріг $1$};
\node[below,font=\small,gray!40!black] at (0,-0.55) {пари проводів};
\end{tikzpicture}
\caption{Виключне АБО на двопровідному коді, лише з \nt{GLUT}-нейронів. Число в кружечку --- поріг в одиницях ваги одного входу. Чотири нейрони І впізнають чотири комбінації входів, два нейрони АБО збирають їх у пару проводів відповіді.}
\label{fig:dualxor}
\end{figure}

\section{Затримки замість годинника}

Для обчислень із часом досить затримок у проводах і детекторів збігів. Найкращий приклад --- те, як мозок визначає, звідки прийшов звук.

Звук від джерела, розташованого збоку, надходить у ближнє вухо раніше, ніж у дальнє. Різниця залежить від напрямку: від нуля, коли джерело прямо попереду, до $0.22\,\text{м}/343\,\text{м/с}\approx0.6\ms$ у людини, коли воно точно збоку. Мозок розрізняє різницю близько десяти \emph{мікросекунд}~\cite{KlumppEady1956,Grothe2010}, хоча самі нейрони спрацьовують із точністю не кращою за частки мілісекунди. Як?

Проведімо від лівого вуха провід уздовж ряду детекторів збігів зліва направо, а від правого --- справа наліво (рис.~\ref{fig:itd}). Сигнал біжить проводом зі швидкістю $v$. До детектора, що стоїть на відстані $x$ від лівого краю ряду довжини $L$, лівий сигнал іде час $x/v$, правий --- $(L-x)/v$. Якщо в праве вухо звук прийшов пізніше на $\delta t$, два сигнали зустрінуться одночасно в тій точці, де
\begin{empheq}[box=\keybox]{equation}
\frac{x}{v} \;=\; \frac{L-x}{v} + \delta t
\quad\Longrightarrow\quad
x \;=\; \frac{L}{2} + \frac{v\,\delta t}{2} .
\label{eq:itd}
\end{empheq}
Спрацює лише той детектор, до якого обидва сигнали прийшли в межах вікна~\eqref{eq:window}. Номер детектора, що спрацював, і є напрямок на джерело. Різниця в часі перетворилася на \emph{місце}.

\begin{figure}[t]
\centering
\begin{tikzpicture}[>=Stealth,
  nrn/.style={circle,draw,very thick,fill=blue!6,minimum size=0.7cm,inner sep=0pt}]
\foreach \k in {0,...,6}{\node[nrn] (D\k) at (1.4*\k,0) {\scriptsize $2$};}
\node[nrn,fill=red!15] at (D4) {\scriptsize $2$};
\draw[->,thick,blue!60!black] (-1.4,0.9) node[left,black] {\small ліве вухо} -- (9.2,0.9);
\draw[->,thick,orange!85!black] (9.8,-0.9) node[right,black] {\small праве вухо} -- (-0.8,-0.9);
\foreach \k in {0,...,6}{
  \draw[->,blue!60!black] (1.4*\k,0.9) -- (D\k.north);
  \draw[->,orange!85!black] (1.4*\k,-0.9) -- (D\k.south);}
\draw[->,thick,red!70!black] (D4.north east) ++(0.1,0.05) -- ++(0.5,0.45) node[above right,font=\small,black] {спрацював};
\node[font=\small,gray!40!black] at (4.2,-1.5) {затримка зростає $\longrightarrow$ \hspace{2.2cm} $\longleftarrow$ затримка зростає};
\pic[scale=1.4] at (-2.3,1.75) {speaker};
\node[font=\scriptsize,align=center] at (-2.3,2.35) {джерело ліворуч};
\end{tikzpicture}
\caption{Визначення напрямку на звук. Сигнали від двох вух біжать назустріч один одному; кожен кружечок --- детектор збігів із порогом $2$. Якщо звук прийшов у праве вухо пізніше, сигнали зустрічаються правіше від середини, за формулою~\eqref{eq:itd}. Вихід мережі --- номер нейрона, що спрацював. Тут джерело ліворуч: у ліве вухо звук прийшов раніше.}
\label{fig:itd}
\end{figure}

Тут немає ні годинника, ні гальмування, ні навіть двопровідного коду: мережа не відповідає «так» чи «ні», а вказує на один нейрон із багатьох. Така схема з ліній затримки й детекторів збігів, схоже, справді працює у слуховому стовбурі птахів~\cite{Carr1990}. Мікросекундна точність виникає не з точності кожного нейрона, а з того, що детекторів багато й кожен дивиться на свою затримку. У ссавців ряду таких затримок не знайшли: там ту саму задачу, схоже, розв'язують детектори збігів, налаштовані точно розрахованим у часі гальмуванням від \nt{GLYC}-нейронів~\cite{Brand2002,Grothe2010}. Як гальмування звужує вікно збігу, ми розберемо в розділі~\ref{sec:gaba} на прикладі \nt{GABA}; гліцин гальмує так само, відкриваючи в отримувача хлорні канали.

\section{Пам'ять}

Комп'ютерові потрібна пам'ять. У такій мережі пам'ять --- це активність, яка підтримує сама себе. Візьмімо групу \nt{GLUT}-нейронів, сильно зв'язаних одне з одним, і опишімо її однією частотою $r$ (рис.~\ref{fig:recurrentgroup}а). У рівновазі $r=f(wr+I)$, де $w$ --- сила зворотного зв'язку групи на саму себе, а $I$ --- приплив ззовні. Розв'яжімо це рівняння графічно, перетинаючи криву $f(wr+I)$ з прямою $r$ (рис.~\ref{fig:bistable}).

\begin{figure}[t]
\centering
% (b): tau dr/dt = -r + f(w r + I), f as in fig:bistable, w=1, I=0.6; Euler dt=0.001 tau.
\begin{tikzpicture}[>=Stealth,
  nrn/.style={circle,draw,thick,fill=blue!6,minimum size=0.5cm,inner sep=0pt}]
\begin{scope}
\node[font=\small] at (-0.5,1.55) {(а)};
\foreach \k in {1,...,6}{\node[nrn] (n\k) at ({1.4+1.0*cos(60*\k)},{1.0*sin(60*\k)}) {};}
\foreach \a/\b in {1/2,2/3,3/4,4/5,5/6,6/1,1/4,2/5,3/6}{\draw[<->,blue!60!black] (n\a) -- (n\b);}
\foreach \k in {2,3,4}{\draw[->,thick,green!45!black] ($(n\k)+(-0.95,0)$) -- (n\k);}
\node[font=\small,green!45!black] at (-0.35,-0.45) {$I$};
\node[font=\Large] at (3.1,0) {$\approx$};
\node[nrn,minimum size=0.85cm,font=\small] (R) at (4.35,-0.1) {$r$};
\draw[->,thick,blue!60!black] (R) edge[loop above,looseness=6] node[above,font=\small] {$w$} (R);
\draw[->,thick,green!45!black] (3.55,-0.95) node[left,font=\small] {$I$} -- (R);
\end{scope}
\begin{scope}[xshift=6.3cm,y=2.6cm,x=1.05cm]
\node[font=\small] at (-0.6,1.25) {(б)};
\draw[->,gray!70!black] (0,0) -- (7.3,0) node[below left,font=\small,black] {$t/\tau$};
\draw[->,gray!70!black] (0,0) -- (0,1.12) node[left,font=\small,black] {$r$};
\foreach \t in {0,2,4,6}{\draw[gray!60!black] (\t,-0.02) -- (\t,0.02); \node[below,font=\scriptsize] at (\t,0) {\t};}
\draw[densely dotted,gray!60!black] (0,0.5) -- (7,0.5) node[right,font=\scriptsize,black,align=left] {поріг\\запису};
\fill[green!45!black,opacity=0.25] (1,-0.2) rectangle (2,-0.13);
\fill[gray!60!black,opacity=0.35] (1,-0.29) rectangle (1.5,-0.22);
\draw[very thick,blue!60!black] plot coordinates {(0.0,0.002) (0.1,0.002) (0.2,0.002) (0.3,0.002) (0.4,0.002) (0.5,0.002) (0.6,0.002) (0.7,0.002) (0.8,0.002) (0.9,0.002) (1.0,0.002) (1.1,0.083) (1.2,0.165) (1.3,0.242) (1.4,0.313) (1.5,0.378) (1.6,0.437) (1.7,0.491) (1.8,0.539) (1.9,0.583) (2.0,0.623) (2.1,0.643) (2.2,0.665) (2.3,0.688) (2.4,0.71) (2.5,0.732) (2.6,0.753) (2.7,0.773) (2.8,0.792) (2.9,0.81) (3.0,0.826) (3.2,0.855) (3.4,0.88) (3.6,0.9) (3.8,0.917) (4.0,0.931) (4.4,0.953) (4.8,0.967) (5.2,0.977) (5.6,0.984) (6.0,0.988) (6.5,0.992) (7.0,0.994)};
\draw[very thick,gray!60!black,densely dashed] plot coordinates {(0.0,0.002) (0.1,0.002) (0.2,0.002) (0.3,0.002) (0.4,0.002) (0.5,0.002) (0.6,0.002) (0.7,0.002) (0.8,0.002) (0.9,0.002) (1.0,0.002) (1.1,0.083) (1.2,0.165) (1.3,0.242) (1.4,0.313) (1.5,0.378) (1.6,0.358) (1.7,0.336) (1.8,0.313) (1.9,0.291) (2.0,0.269) (2.2,0.228) (2.4,0.191) (2.6,0.159) (2.8,0.133) (3.0,0.11) (3.2,0.091) (3.4,0.076) (3.6,0.063) (3.8,0.052) (4.0,0.043) (4.4,0.03) (4.8,0.021) (5.2,0.015) (5.6,0.011) (6.0,0.008) (6.5,0.006) (7.0,0.004)};
\node[font=\scriptsize,blue!60!black,anchor=west] at (3.3,0.8) {приплив $1\tau$: записано};
\node[font=\scriptsize,gray!50!black,anchor=west] at (2.3,0.3) {приплив $0.5\tau$: згасло};
\end{scope}
\end{tikzpicture}
\caption{Група, сильно зв'язана сама з собою. (а)~Багато \nt{GLUT}-нейронів, що збуджують одне одного, поводяться як один нейрон із частотою $r$, який збуджує сам себе із силою $w$; ззовні надходить приплив $I$. (б)~Частота групи в часі при $w=1$ і тій самій кривій $f$, що на рис.~\ref{fig:bistable}. Приплив $I=0.6$ подається на час, позначений смужкою під віссю. Якщо він триває $1\tau$, група встигає перейти нестійку точку $r=0.5$, розгоряється й лишається горіти, коли приплив закінчився. Якщо лише $0.5\tau$, вона не дотягує до цієї точки й згасає назад: щоб записати біт, приплив має тривати довше приблизно $0.7\tau$.}
\label{fig:recurrentgroup}
\end{figure}

\begin{figure}[t]
\centering
\begin{tikzpicture}[>=Stealth,x=5cm,y=3.2cm]
\draw[->,gray!70!black] (0,0) -- (1.1,0) node[below left,black] {\small $r$};
\draw[->,gray!70!black] (0,0) -- (0,1.1);
\draw[thick,gray!60!black] (0,0) -- (1.05,1.05) node[right,black] {\small $r$};
\draw[very thick,blue!60!black] plot[domain=0:1.05,samples=80] (\x,{1/(1+exp(-(\x-0.5)/0.08))});
\draw[very thick,red!70!black,densely dashed] plot[domain=0:1.05,samples=80] (\x,{1/(1+exp(-(\x+0.3-0.5)/0.08))});
\fill[blue!60!black] (0.002,0.002) circle (0.018);
\draw[blue!60!black,fill=white,thick] (0.5,0.5) circle (0.018);
\fill[blue!60!black] (0.998,0.998) circle (0.018);
\node[below right,font=\small] at (0.02,0.0) {вимк.};
\node[right,font=\small] at (0.52,0.46) {нестійко};
\node[below,font=\small] at (0.99,0.96) {увімк.};
\node[anchor=east,font=\small,red!70!black] at (0.19,0.62) {$f(wr+I)$};
\node[anchor=west,font=\small,blue!60!black] at (0.45,0.1) {$f(wr)$};
\end{tikzpicture}
\caption{Пам'ять із групи \nt{GLUT}-нейронів із сильним зворотним зв'язком. Суцільна крива --- без зовнішнього припливу: два стійкі перетини з прямою $r$, «вимкнено» й «увімкнено», і нестійкий між ними. Короткий приплив $I$ (штрихова крива) залишає лише верхній.}
\label{fig:bistable}
\end{figure}

Якщо зворотний зв'язок сильний, перетинів три: нижній --- тиша, верхній --- група горить на повну, середній нестійкий. Вийшов елемент із двома стійкими станами, тригер. Записати в нього одиницю легко: короткий приплив піднімає криву, нижній перетин зникає, і група розгоряється --- і лишається горіти, коли приплив закінчився (рис.~\ref{fig:recurrentgroup}б). Надто короткий приплив не встигає довести групу до нестійкої точки, і вона згасає назад. Таку самопідтримувану активність, що триває секунди після того, як подразник зник, справді спостерігають у мозку й вважають основою короткочасної пам'яті~\cite{Wang2001}. Але це не єдиний спосіб пам'ятати секунди. Запис можна зберігати й у повільній змінній самих синапсів: після пачки імпульсів полегшення, як у синапсу-«тире» розділу~\ref{sec:code}, триває близько секунди, і мережа між короткими спалахами може майже мовчати --- не тригер, а заряджений конденсатор~\cite{Mongillo2008}. Такі «тихі» проміжки під час утримання в пам'яті спостерігають~\cite{Stokes2015}, і зберігання без постійних імпульсів дешевше за енергією. Який із двох способів кора використовує в якому випадку, обговорюється.

А як стерти? За твердженням~\ref{prop:monotone} активністю --- ніяк; лишається щось \emph{прибрати}. Перший спосіб --- прибрати підтримку: якщо група горить лише разом із припливом від іншої групи, пам'ять згасає разом із ним. Другий --- утома. У справжніх синапсів за тривалої роботи запас медіатора виснажується~\cite{Tsodyks1997}, а в нейронів наростають повільні струми, що знижують збудливість. Зворотний зв'язок слабшає, верхній перетин зникає, і група згасає сама через секунди. Виходить пам'ять, яка вмикається сигналом, а вимикається лише часом.

\section{Послідовності}

Замість годинника можна мати \emph{таймер, що запускається подією}. З'єднаймо групи \nt{GLUT}-нейронів у ланцюжок: перша збуджує другу, друга --- третю, і так далі. Зовнішній сигнал запалює першу групу, і хвиля активності пробігає ланцюжком, кожна ланка через одну-дві затримки після попередньої і лише один раз: відразу після імпульсу нейрони несприйнятливі, а хвиля вже пішла далі.

Такий ланцюжок відмірює час від події: спрацювала ланка номер $k$ --- отже, від моменту запуску минуло $k$ затримок. Його виходи можна подати на м'язи й отримати послідовність рухів, яка розгортається сама. У співочих птахів під час співу окремі \nt{GLUT}-нейрони одного з відділів мозку спрацьовують суворо по черзі, кожен у свій момент пісні, --- дуже схоже на такий ланцюжок~\cite{Hahnloser2002}.

На відміну від годинника, ланцюжок мовчить, доки його не запустили, відпрацьовує один раз і відмірює час лише для тих, хто до нього під'єднаний. Спільного такту в мережі й далі немає.

\section{Чого бракує}

З самих лише \nt{GLUT}-нейронів, без гальмування, виходить чимало. Але трьох речей бракує, і всіх трьох через твердження~\ref{prop:monotone}.

\emph{Вибір.} Мережа має обирати один напрямок, одну дію. Якщо два подразники майже рівні, у мережі рис.~\ref{fig:itd} спрацюють обидва виходи, і придушити слабший на користь сильнішого нічим.

\emph{Скидання.} Стерти пам'ять сигналом не можна.

\emph{Стійкість.} У мережі, де зв'язки виникають не за планом інженера, петлі додатного зворотного зв'язку неминучі, й активність у них може наростати лавиною (розділ~\ref{sec:neurontypes}). Єдине гальмо --- та сама втома, повільна й груба.

Усі три біди лікуються одними ліками --- нейроном, що гасить імпульс імпульсом. Це \nt{GABA}, і потрібен він не для того, щоб обчислювати, а для того, щоб обирати, скидати й тримати обчислення під контролем.

\section{Задачі}

\begin{exercise}[\lvlA]
\label{ex:02:1}
\emph{Ряд детекторів для звуку.} Сигнали від вух біжать проводами рис.~\ref{fig:itd} зі швидкістю $5$~м/с, найбільша різниця приходу $0.64\ms$. Детектори --- нейрони~\eqref{eq:glutneuron} з $\tau=0.5\ms$, $\theta=1.5w$ --- стоять так, що сусідні розрізняють $10$~мкс. Скільки детекторів у ряду і скільки з них спрацює на один звук?
\end{exercise}

\begin{exercise}[\lvlB]
\label{ex:02:2}
\emph{Нерівні входи зсувають вікно.} Детектор збігів~\eqref{eq:glutneuron} з $\tau=0.5\ms$, $\theta=1.5$ отримує по імпульсу від двох входів із вагами $1$ і $0.8$. Знайдіть вікно збігу та його центр. На скільки мікросекунд помилиться ряд рис.~\ref{fig:itd}, якщо вхід від одного вуха на $20\,\%$ слабший, ніж від іншого?
\end{exercise}

\begin{exercise}[\lvlB]
\label{ex:02:3}
\emph{Які мережі не коливаються.} Мережа $\tau\dot r_i=-r_i+f_i\bigl(\textstyle\sum_jw_{ij}r_j+I_i\bigr)$ з неспадними $f_i$ і $w_{ij}\ge0$ при $i\ne j$ монотонна й стійко не коливається. Доведіть: заміною $r_i\to\pm r_i$ до неї зводиться мережа з парною кількістю гальмівних зв'язків у кожному циклі. Чи може коливатися взаємне гальмування двох груп? Петля \nt{GLUT}--\nt{GABA}?
\end{exercise}

\begin{exercise}[\lvlB]
\label{ex:02:4}
\emph{Проєктування: двопровідний суматор.} Кожен біт передається парою проводів $(x,\bar x)$. Побудуйте повний суматор трьох бітів, що видає пари $(S,\bar S)$ суми та $(C,\bar C)$ перенесення, з \nt{GLUT}-нейронів, указавши ваги й пороги. Вкладіться у вісім нейронів. Скільки потребувала б схема з елементів І та АБО, як на рис.~\ref{fig:dualxor}?
\end{exercise}

\begin{exercise}[\lvlB]
\label{ex:02:5}
\emph{Моделювання: скільки триває запис.} Група рис.~\ref{fig:recurrentgroup}: $\tau\dot r=-r+f(r+I)$, $f(u)=1/\bigl(1+e^{-(u-0.5)/0.08}\bigr)$, до запису --- у нижньому стані; приплив $I$ подається на час $T$. Знайдіть чисельно найменше $T$, що записує біт при $I=0.6$, і поріг $I_c$, нижче якого біт не записується за жодного $T$.
\end{exercise}

\begin{exercise}[\lvlA]
\label{ex:02:6}
\emph{Ланцюжок як таймер.} Ланка ланцюжка --- $100$ \nt{GLUT}-нейронів, затримка $2\ms$ з незалежним розкидом $0.2\ms$. (а)~Скільки нейронів потрібно, щоб відміряти $1$~с, і яка відносна похибка? Як вона залежить від інтервалу? (б)~У справжнього таймера похибка --- $5\,\%$ за будь-якого інтервалу. Яке джерело розкиду це дає?
\end{exercise}

\begin{exercise}[\lvlC]
\label{ex:02:7}
\emph{Дослідження: тригер чи конденсатор.} $100$ нейронів тримають біт $10$~с. Тригер: кожен дає $20$ імпульсів за секунду. Конденсатор: полегшення $u$ спадає з $\tau_F=1$~с, біт читається при $u\ge u_{\max}/2$, освіження --- пачка з трьох імпульсів на нейрон. У скільки разів конденсатор ощадливіший і що буде з бітом після глушіння групи на $200\ms$?
\end{exercise}
